\chapter*{Introduction Générale}
\label{chapter:introduction}
\addcontentsline{toc}{chapter}{Introduction}
\markboth{Introduction Générale}{}

Dans le cadre de la formation d'ingénieur en Informatique et Réseaux à l'École Marocaine
des Sciences de l'Ingénieur (EMSI), les étudiants doivent effectuer un stage de
perfectionnement en entreprise afin de mettre en pratique leurs connaissances théoriques
dans un contexte professionnel réel. C'est dans ce cadre que ce stage de deux mois s'est
déroulé au sein du département Informatique (IT) du Mövenpick Hôtel de Tanger.

Un réseau local d'hôtel voit se connecter, au fil du temps, un grand nombre
d'équipements hétérogènes : postes de travail, imprimantes, caméras, bornes Wi-Fi et
appareils invités. Sans outil dédié, le département IT ne dispose d'aucun moyen simple
de savoir, à un instant donné, combien d'appareils sont actifs sur le réseau, ni
d'identifier une connexion inhabituelle ou non autorisée. C'est pour répondre à ce
besoin qu'est né le projet \textbf{NetSentry} : un outil interne de découverte et de
cartographie des équipements connectés au réseau local, capable de scanner le réseau, de
conserver un historique des détections et de signaler automatiquement tout équipement
inconnu.

Le présent document constitue le rapport de stage du projet NetSentry. Il a pour objectif
de présenter, de manière claire, structurée et détaillée, le contexte du stage, l'étude
de l'existant et la méthodologie de travail suivie, la conception du système, son
architecture technique et les technologies mobilisées, ainsi que les résultats obtenus.

Ce rapport s'organise en cinq chapitres.

Le \textbf{premier chapitre} présente le contexte général du projet : l'entreprise
d'accueil, le département Informatique, la problématique à l'origine du projet ainsi que
les objectifs poursuivis.

Le \textbf{deuxième chapitre} est consacré à l'étude de l'existant et à la méthodologie
de travail : les solutions de supervision réseau déjà disponibles, leurs limites, le
cahier des charges qui a cadré le stage, ainsi que les exigences fonctionnelles et non
fonctionnelles retenues.

Le \textbf{troisième chapitre} présente l'analyse et la conception du système :
diagramme de cas d'utilisation, diagramme de séquence et modèle conceptuel de données.

Le \textbf{quatrième chapitre} détaille l'architecture technique et les technologies
mobilisées -- avec une attention particulière portée à Nmap, l'outil au c\oe{}ur du moteur de
scan -- ainsi que l'organisation du travail retenue pour la réalisation du projet
(planning prévisionnel et suivi des tâches).

Le \textbf{cinquième chapitre} présente, à travers des captures d'écran, la mise en
\oe{}uvre concrète de l'outil développé : tableau de bord, détail d'un appareil, alertes,
historique des scans, ainsi que l'historique de versionnage du projet sur GitHub.
